% !TeX root = ../papers/nc-p3.tex
\section{二つの球面twistによる楕円繊維間の移送}
\label{sec:finite-elliptic-fiber-transfer}
\begin{lemma}\label{lem:octic-linear-determinant}
前節のいずれかの滑らかな楕円八次繊維を $F$、$L=\mathcal O_F(h)$、
$D_0=H^0(\mathbb P^3,\mathcal O(1))\subset H^0(L)$ とする。
全ての次数4の線束 $J$ で $D_0\otimes H^0(J)\to H^0(LJ)$ は全射、核 $K_J$ は次元4である。
一般の $x\in T$ と一般の $J$ で、この核を $x$ で評価する写像
$K_J\to H^0(J)\otimes\mathcal O(h)|_x$ は同型となる。
\end{lemma}
\begin{proof}
$I_F(4)$ は大域生成されるので $F$ と長さ5以上で交わる直線はない。
実際その直線を全ての四次切断が含めば、基底軌跡に含まれてしまう。
$M=\ker(D_0\otimes\mathcal O_F\to L)$ とする。
次数 $d\le3$ の線束 $J_d$ による消える積
$t\in D_0\otimes H^0(J_d)$ の非零元はテンソル階数3以上である。
階数1は整性に反する。階数2なら係数の二切断の固定因子を除いた鉛筆は次数 $e\le d$ であり、
対応する二つの超平面切断はこの鉛筆に共通の切断を掛けたものになる。
従って超平面の交線は $F$ と長さ $8-e\ge5$ で交わる。
特に $M^*$ は次数2以下の商直線束を持たない。実際その商への
$D_0^*\otimes\mathcal O_F$ の写像は上の形の関係であり、次数2以下では切断が高々二つである。

$M^*$ には次数4以上の部分直線束もない。飽和してそのような束を割ると階数2・次数4以下の商 $Q$ が残る。
任意のそのような $Q$ は次数2以下の商直線束を持つ。さもなくば全ての次数2の $B$ で
$H^0(Q^* B)=0$ であるが、点 $p$ を一つ選べばRiemann--Rochにより
$h^0(Q^*B(p))\ge2$、これを各点 $q$ で評価した核は $H^0(Q^*B(p-q))=0$。
従って切断の評価は全点で同型となるが行列式の次数が正で矛盾する。
さらに次数3の部分直線束は高々一つである。二本が異なる像を持てばそれらの和の飽和を割って
次数2以下の商を得るからである。同じ像に対する写像は定数倍しかない。
従ってSerre双対より全ての次数4で $H^1(MJ)=0$、これが全射性と核次元を示す。
次数3については $h^0(MJ_3)=1$ で、ただ一つあり得る例外でも高々2である。

$x\notin F$ で $D_x\subset D_0$ を $x$ で消える三切断とし、
$M_x=\ker(D_x\otimes\mathcal O_F\to L)$ とする。この評価は基点を持たない。
$M_x^*$ が半安定でなければ次数3以下の商直線束、従って $D_x$ に支えられる上記の関係がある。
その階数は3、従って次数3で、関係から $D_x$ は一意に決まる。
次数3のPicard曲線の各点で関係の射影空間の次元は零、一つの例外でも高々1であるから、
このような $x$ の族は次元1以下である。これは点の個数からの推測ではなく、
Picard曲線上で切断の積の核を射影化した有限型の入射対応の次元評価である。
従って一般の $x\in T$ で $M_x^*$ は半安定である。
階数2・傾き4の半安定束に入る次数4の線束の同型類は高々二つである。
一本を入れると商も同次数の線束で、他のものはこの部分または商に同型だからである。
よって一般の $J$ に対して $H^0(M_x J)=0$、すなわち $K_J$ の $x$ 評価が単射である。
\end{proof}

\begin{theorem}\label{thm:octic-fiber-twist-transfer}
滑らかな $F\in|F|$ に次数 $-4$ の一般の線束 $J^{-1}$ を置き、その直像を $X$ とする。
$A_0=T_{\mathcal O}T_{\mathcal O(h)}$ に対し、$A_0^{-1}X=Y[-1]$ である。
$Y$ は $|C|$ の一繊維にスキームとして支えられる純一次元層で、Mukaiベクトルは $(0,C,12)$ である。
従って $\Phi[-2](X)\simeq X$ である。$|C|$ の全繊維が整なら $Y$ はその繊維上で階数1の無捩れ層である。
\end{theorem}
\begin{proof}
$i:F\hookrightarrow T$ とする。$J$ は四切断で生成され、評価の核
$W=\ker(H^0(J)\otimes\mathcal O_T\to i_*J)$ は局所自由である。
実際滑らかなCartier因子上の線束は局所射影次元1である。
$\mathcal O_F(F)$ は自明なので、双対列は
$0\to H^0(J)^*\otimes\mathcal O_T\to W^*\to X\to0$ である。
これは $T_{\mathcal O}^{-1}X$ を与える。逆twistの評価は
$\operatorname{Ext}^1_T(X,\mathcal O_T)=H^0(J)$ の恒等テンソルによる普遍拡大であり、
上の双対列は評価そのものを双対化しているからである。
次に $H^*(T,W(h))$ は前補題の積の核 $K_J$ が次数零にあるだけである。
従って $T_{\mathcal O(h)}^{-1}W^*$ は $W^*\to K_J^*\otimes\mathcal O(h)$ の錐の $-1$ 移動である。
$F$ の外ではこれは前補題の四行四列評価の転置で、一般の $J$ で一般点同型である。
束間の単射なので余核 $Y$ は純一次元である（局所射影次元1）。
行列式は $4h-F=C$ の非零切断であり、その随伴行列によって行列式の方程式が余核を消す。
鉛筆は連結繊維なのでこの零因子は実際に一つの繊維で、その上で $\mathcal O_T(C)$ は自明である。
ベクトルは命題\ref{prop:quartic-k3-two-pencil-twists}の計算で決まる。
整な場合の階数は一次の重複度が1であることから従う。
$X$ 上での $L_{-F}$ の自明性と $Y$ 上での $L_{-C}$ の自明性を分解式へ代入すれば最後の同型を得る。
\end{proof}
